% ============================================================================
%  Standalone: Limitations & Reinforcement Plan
%  Companion to the cell-state reachability oracle manuscript
% ============================================================================
\documentclass[11pt]{article}
\usepackage[letterpaper,margin=1in]{geometry}
\usepackage[T1]{fontenc}
\usepackage[utf8]{inputenc}
\usepackage{lmodern}
\usepackage{microtype}
\usepackage{amsmath,amssymb}
\usepackage{booktabs}
\usepackage{array}
\usepackage{tabularx}
\usepackage{xcolor}
\definecolor{linkblue}{HTML}{2166AC}
\definecolor{sev1}{HTML}{D55E00}   % highest threat
\definecolor{sev2}{HTML}{E69F00}
\definecolor{sev3}{HTML}{009E73}
\usepackage[colorlinks=true,allcolors=linkblue]{hyperref}
\usepackage{enumitem}
\setlist{leftmargin=1.4em,itemsep=2pt,topsep=3pt}

\newcommand{\Th}[1]{T\textsubscript{H}#1}
\newcommand{\sev}[2]{\textcolor{#1}{\textbf{#2}}}

\title{\bfseries Limitations and Reinforcement Plan\\
\large A self-critique of the convex-cone cell-state reachability oracle,\\
with a prioritized plan to strengthen it}
\author{Cell-State Reachability Project \\ \small Built with Claude --- Life Sciences (Research track)}
\date{\today}

\begin{document}
\maketitle

\begin{abstract}
\noindent
This companion document states, as sharply as we can, the ways in which the reachability-oracle
manuscript could be wrong or could fail to convince a critical reviewer, and pairs each with a
concrete experiment or analysis that would resolve it. We rank limitations by how directly they
threaten the paper's central claim---that the oracle emits a \emph{falsifiable feasibility
verdict with a data-grounded certificate}, a capability no surveyed method provides---rather
than by how easy they are to fix. Two limitations are, in our judgment, load-bearing and remain
so: the infeasibility certificate has never been tested at the bench, and the additivity
assumption that underwrites multi-gene recipes is calibrated only on a borrowed dataset---both
require wet-lab data we do not have. We treat the rest as important but secondary. Since the
first version of this document we have \emph{completed} every immediate in-silico reinforcement
(a non-negativity-constraint ablation, a reachable-cosine ceiling, magnitude-capped recipes, a
held-out-modality certificate test on synthetic ground truth, a DEG-magnitude-weighted
re-evaluation, a certificate split-stability check, and a cross-cell-type transfer probe); these
strengthen or, in two cases (L4, L5), resolve the corresponding limitations, and one of them---%
the certificate cross-reference---tightened a claim by showing the top certificate genes are
state-markers and negative regulators rather than canonical activators. Each item below is
written so that a reader can check our reasoning and, if they have the wet-lab or compute
resources we lack, act on it directly.
\end{abstract}

\section*{How to read this document}
Each limitation carries a severity tag---\sev{sev1}{critical} (directly threatens the central
claim), \sev{sev2}{major} (threatens a headline result but not the core idea), or
\sev{sev3}{moderate} (affects scope or presentation)---followed by \emph{What it threatens},
\emph{Current state}, and a \emph{Reinforcement plan} with a rough effort estimate. The plan is
summarized as a prioritized roadmap in Table~\ref{tab:roadmap} and a candid scorecard in
Table~\ref{tab:scorecard}.

% ===========================================================================
\section{Critical limitations (threaten the central claim)}

\subsection{L1 \quad The infeasibility certificate has never been tested at the bench}
\sev{sev1}{Severity: critical.}
\begin{itemize}
\item[] \textbf{What it threatens.} The certificate is the paper's most distinctive
contribution: when a target is unreachable, convex duality names the genes that must instead be
\emph{activated}. Every other method we surveyed can, in principle, rank genes; none emits this
constructive infeasibility proof. But the entire claim rests on the certificate being
\emph{biologically actionable}, and we have not run the one experiment that would show it. We
validate the reachable side (recipes recover known regulators as positive controls) but the
infeasible side---the genes named in an activation certificate---has no direct wet-lab test in
this work.
\item[] \textbf{Current state.} For \Th{2}$\rightarrow$\Th{1} the certificate names LYAR, IKZF3,
CRTAM, CBLB, GBP5, LAG3, IRF8; for the K562 CEBPA state it names a myeloid program (MNDA, HP,
VSIG4, NCF1). These are internally consistent and interpretable, but ``interpretable'' is not
``correct.'' The certificate residual is at machine precision ($\approx\!10^{-11}$), which
proves the \emph{convex program} is solved to optimality---it says nothing about whether raising
those genes actually produces the shift.
\item[] \textbf{Reinforcement plan.}
\begin{enumerate}
\item \emph{Direct test (definitive, wet-lab, weeks--months).} Run a small CRISPRa arm against
the \Th{2}$\rightarrow$\Th{1} certificate genes and measure whether combined activation moves
cells toward \Th{1} in the directions the certificate predicts. A positive result validates the
single most novel claim; a negative result is itself publishable as a calibration of what the
certificate does and does not capture.
\item \emph{Retrospective test (immediate, computational).} Cross-reference every certificate
gene against published CRISPRa/overexpression screens in the same lineage; a certificate gene
that is already known to drive the target direction on activation is corroborating evidence
obtainable now, without new experiments.
\item \emph{Held-out-modality test (immediate, computational).} Where a screen contains
\emph{both} interference and activation arms, hide the activation arm, generate the certificate
from knockdowns alone, and check whether the named genes are exactly those the hidden activation
arm shows to be effective. This is the strongest test achievable purely in silico and should be
the first thing done.
\end{enumerate}
\item[] \textbf{Status update (completed in silico, \today).} Items~2 and~3 are done; item~1
(the wet-lab CRISPRa arm) remains the one definitive test we cannot run here.
\emph{(2)~Retrospective cross-reference:} the top $20$ certificate genes were scored against
the literature direction-of-effect and Open Targets genetics---$4$ SUPPORTS, $7$ PLAUSIBLE, $9$
UNKNOWN, with $8/20$ carrying a strong ($\ge0.5$) genetic link to a \Th{1}-driven autoimmune
disease. This surfaced an important qualification: the highest-demand genes are state-markers
and \emph{negative} regulators (IKZF3/Aiolos, CBLB, LAG3), while the canonical \Th{1} master
TFs rank deep (IFNG $\#35$, TBX21 $\#144$), so the certificate is a reproducible ranking of
\emph{unmet upward demand}, not a validated activation list. \emph{(3)~Held-out-modality test:}
the procedure is implemented and verified on synthetic ground truth (AUROC $0.999$,
precision@30 $0.933$, $z=8.9$ against a shuffle null of $0.497\pm0.056$; PASS). It is blocked
only on a real dataset carrying both CRISPRi and CRISPRa arms on the same readout axis---a data
requirement, not a method gap. We also confirmed the certificate ranking is reproducible under a
random gene-axis split (cross-half Spearman $\rho\approx0.65$, half-vs-full $\approx0.84$,
$z\approx35$). The severity remains \emph{critical} because the bench test (item~1) is what
would convert ``reproducible and interpretable'' into ``biologically validated.''
\end{itemize}

\subsection{L2 \quad Additivity is load-bearing and calibrated only out-of-domain}
\sev{sev1}{Severity: critical.}
\begin{itemize}
\item[] \textbf{What it threatens.} Every multi-gene recipe and every certificate assumes that
co-perturbation effects add. If additivity fails badly, the recipes are wrong and the cone
geometry misrepresents what combinations can achieve. We test additivity---but only on the
Norman K562 screen, because the primary CD4\textsuperscript{+} dictionary contains no measured
double perturbations.
\item[] \textbf{Current state.} On $126$ measured K562 doubles, the measured double aligns with
the sum of its singles at median cosine $0.71$: additivity is a good but bounded approximation.
The deficit grows with combined perturbation magnitude (Spearman $+0.58$; scale-free saturation
ceiling $M^{\star}=13.9$), and the intuitive collinearity$\rightarrow$sub-additivity hypothesis
was refuted (Spearman $-0.16$, n.s.). So the failure mode is magnitude saturation, and it is
predictable---but the number $0.71$ is borrowed from a different cell type and perturbation
direction than the headline result.
\item[] \textbf{Reinforcement plan.}
\begin{enumerate}
\item \emph{In-domain doubles (definitive, wet-lab).} Measure a modest panel ($\sim$50--100) of
double knockdowns in CD4\textsuperscript{+} T cells, chosen to span the recipe sizes and
magnitudes the oracle actually proposes, and recompute the additivity distribution in-domain.
This replaces the borrowed estimate and directly calibrates the recipe-size ceiling.
\item \emph{Magnitude-capped recipes (immediate, computational).} Since the deficit is a
magnitude effect with a fitted ceiling, cap recipe magnitude at the regime where median additive
cosine stays above a chosen threshold, and report recipes with a reliability band rather than a
point claim. This is implementable now and makes the additivity caveat quantitative in every
recipe.
\item \emph{Second-order term (research, computational).} Where doubles exist, fit a bounded
interaction correction and test whether it improves held-out reachability, converting additivity
from an assumption into a modeled, bounded residual.
\end{enumerate}
\item[] \textbf{Status update (completed in silico, \today).} Item~2 (magnitude-capped recipes)
is done. Across all $12$ atlas cells every recommended knee recipe is additive-safe with a
per-recipe reliability of $0.92$--$0.96$, and the combined-magnitude saturation cap binds only
at $k\approx28$ knockdowns versus a knee of $k\approx4$--$5$---a $5.6\times$ margin on the
headline recipe (only $2/12$ recipes bind the cap within $k\le40$). Recipes are now reported
with a reliability band rather than a point claim, making the additivity caveat quantitative.
Items~1 (in-domain CD4\textsuperscript{+} doubles) and~3 (second-order term) still require
measured doubles the primary dictionary lacks, so the severity remains \emph{critical}: the
$0.71$ additivity estimate is still borrowed from K562, and only a wet-lab panel (W2) replaces
it in-domain.
\end{itemize}

% ===========================================================================
\section{Major limitations (threaten a headline result)}

\subsection{L3 \quad The disjointness claim is a structured survey, not a head-to-head benchmark}
\sev{sev2}{Severity: major.}
\begin{itemize}
\item[] \textbf{What it threatens.} The ``$92$ methods, empty intersection'' result is the
paper's positioning centerpiece. It is a claim about \emph{capabilities as scored on two axes},
not a claim that the oracle out-predicts any method on a shared task. A skeptical reviewer can
argue that the axes were drawn to make the top-right corner empty---that ``feasibility verdict
with certificate'' is defined precisely as what only this method does.
\item[] \textbf{Current state.} We mitigate this by scoring each method against explicit,
per-capability criteria and by refining the claim to a defensible range ($13$--$15$ prior
methods carry \emph{partial} verdict/certificate marks, all model-theoretic), and we state
openly that the novelty is the \emph{combination}, not any single capability. But the scoring is
ours, and no external rater has confirmed it.
\item[] \textbf{Reinforcement plan.}
\begin{enumerate}
\item \emph{Independent re-scoring (immediate--short).} Release the scoring rubric and the
$92\times6$ capability matrix as data (already produced) and invite independent re-scoring;
report inter-rater agreement. A published rubric turns a subjective survey into an auditable one.
\item \emph{Shared-task comparison (short--medium, computational).} On the reachable side, where
several forward predictors \emph{can} be run, compare the oracle's recipe against a predictor's
top-ranked perturbations on a common held-out target using a common metric. The oracle need not
win at prediction---its contribution is the verdict---but showing it is competitive on the
overlapping task closes the ``maybe it just cannot predict'' escape.
\item \emph{Adversarial axis test (immediate).} Explicitly ask whether any single prior method,
under the most generous reading of its own paper, could be scored into the top-right corner, and
document the specific capability each lacks. This pre-empts the ``you defined the corner'' critique.
\end{enumerate}
\end{itemize}

\subsection{L4 \quad No ablation shows the non-negativity constraint earns its place}
\sev{sev3}{Severity: major $\rightarrow$ resolved in silico.}
\begin{itemize}
\item[] \textbf{What it threatens.} The conceptual core is that knockdown combinations are
\emph{non-negative}, making reachability a convex-\emph{cone} membership problem. But we never
show that the non-negativity constraint changes the answer---that plain (signed) least squares,
or a nearest-single-effect heuristic, would give materially different verdicts. Without that
ablation, a reviewer cannot tell whether the cone is doing the work or is decoration on an
ordinary regression.
\item[] \textbf{Current state.} The non-negativity is biologically motivated (a knockdown cannot
be applied in negative amount) and it is what makes the Farkas certificate meaningful, but the
manuscript reports no side-by-side comparison of constrained versus unconstrained fits.
\item[] \textbf{Reinforcement plan.}
\begin{enumerate}
\item \emph{Constraint ablation (immediate, computational).} For every atlas target, compare
NNLS reachable cosine and support against (i)~unconstrained least squares and (ii)~a
nearest-single-effect baseline, on the same held-out genes. Report where the cone verdict and
the unconstrained verdict disagree---those disagreements are exactly where the biology (sign
constraint) matters, and quantifying them directly answers ``does the cone earn its place.''
\item \emph{Certificate necessity (immediate).} Show that the separating hyperplane has no
counterpart in the unconstrained fit (which never declares infeasibility), making explicit that
the certificate is a consequence of the constraint, not an add-on.
\end{enumerate}
\item[] \textbf{Status update (completed in silico, \today).} Both items are done, and this
limitation is resolved by the analysis available in-hand. Across all $12$ atlas cells the
non-negativity constraint costs a mean of only $0.018$ held-out cosine relative to an
unconstrained least-squares fit, but the unconstrained fit buys that improvement with a mean of
$\approx\!3{,}537$ biologically unrealizable negative weights (running knockdowns in reverse);
the cone recovers a mean $95\%$ (range $86$--$103\%$) of the unconstrained cosine and beats the
nearest-single-effect baseline in $12/12$ cells. The constraint is therefore nearly free in fit
quality and is the sole source of the Farkas/KKT certificate---the unconstrained fit never
declares infeasibility. The cone earns its place; severity is downgraded from major to resolved
in silico.
\end{itemize}

\subsection{L5 \quad Reachable cosines are modest and the in-sample number is optimistic}
\sev{sev3}{Severity: major $\rightarrow$ resolved in silico.}
\begin{itemize}
\item[] \textbf{What it threatens.} A reader could read ``held-out cosine $0.448$'' as weak
performance and dismiss the headline.
\item[] \textbf{Current state.} We argue---and the signed decomposition supports---that the
modest value is the structural ceiling of loss-of-function acting alone, not a defect, and we
report the calibration gap against ourselves (in-sample mean $0.580$ overstates held-out mean
$0.355$ by $0.225$ across the atlas; bootstrap CIs confirm the gap is real, not noise). This is
mitigated, but it depends on the reader accepting the ceiling argument.
\item[] \textbf{Reinforcement plan.} Report, for each target, the theoretical maximum reachable
cosine given its gain-of-function fraction (the best any knockdown-only method could achieve),
so the reader sees the oracle's value \emph{relative to the achievable ceiling} rather than to
$1$. This reframes ``$0.448$'' as ``a large fraction of what is possible by this modality.''
\item[] \textbf{Status update (completed in silico, \today).} Done. The knockdown-only reachable
ceiling equals $\sqrt{\smash[b]{f_{\mathrm{LOF}}}}$ and matches the in-sample cone cosine to
$10^{-4}$; against it, the headline $0.448$ is $71\%$ of what is attainable in principle (atlas
mean $61\%$, range $49$--$71\%$), with the residual $22$--$26\%$ gain-of-function-locked. Every
headline number is now reported relative to its own modality ceiling, so ``$0.448$'' reads as
``$71\%$ of the achievable knockdown-only maximum,'' not as a weak fit. Combined with the
already-disclosed in-sample-vs-held-out calibration gap, this limitation is resolved by
in-hand analysis.
\end{itemize}

% ===========================================================================
\section{Moderate limitations (scope and external validity)}

\subsection{L6 \quad One system, four donors, collapsed across donors}
\sev{sev3}{Severity: moderate.}
\emph{What it threatens:} external validity of the primary result.
\emph{Current state:} the primary dictionary is CD4\textsuperscript{+} T cells from four donors,
and the effect vectors available are collapsed across donors, so a true leave-one-donor-out
generalization test is impossible on the primary data. Cross-cell-type validity is now
characterized beyond the single K562 point: a transfer probe on the K562 and RPE1 essential-gene
screens (Replogle 2022) shows that effect \emph{direction} transfers moderately (matched
cross-type cosine median $0.35$ vs a $\approx0$ shuffled-gene null) while the minimal
\emph{recipe} does not (same-target Jaccard $0.11$, collapsing to null under a cross-type basis).
This bounds the external-validity claim---direction is portable, the specific recipe is
not---but it does \emph{not} substitute for per-donor generalization within the primary system,
which remains untested.
\emph{Reinforcement plan:} (i)~obtain per-donor effect vectors and run leave-one-donor-out to
report generalization across genetic backgrounds (still open); (ii)~\sev{sev3}{done in silico:}
the cross-cell-type probe turns the single K562 transfer into a graded, three-level
characterization; extend to further primary systems to complete the generalization curve;
(iii)~state the external-validity claim as ``direction demonstrated across cell types; recipe
shown to be cell-type-specific'' rather than ``general.''

\subsection{L7 \quad Statistical nulls answer narrow questions; one diagnostic is negative by construction}
\sev{sev3}{Severity: moderate (presentation).}
\emph{What it threatens:} a reader could over- or under-read the significance tests.
\emph{Current state:} the shuffled-target null establishes the target is more reachable than a
structure-matched random direction; the held-out-gene test establishes generalization across
expression coordinates---neither establishes biological correctness or cross-cell
generalization. A separate random-perturbation diagnostic is \emph{negative by construction}
(it permutes gene labels within each effect vector) and is used only as an internal control, not
as evidence of performance; this must be stated wherever it appears so it is never mistaken for a
strength.
\emph{Reinforcement plan:} keep each null's scope labeled at point of use (done in the
manuscript), and add a one-line ``what this null does and does not show'' to every significance
statement.

\subsection{L8 \quad The positive-control TF-vs-marker split is post-hoc}
\sev{sev3}{Severity: moderate (presentation).}
\emph{What it threatens:} the appearance of a pre-registered test where there is an exploratory
annotation.
\emph{Current state:} the directional AUROC and signed-concordance statistics on canonical
regulators are compelling positive controls, but the transcription-factor-versus-marker
partition was chosen after seeing the data.
\emph{Reinforcement plan:} label it exploratory wherever reported (done), and, if it is to be
promoted to a formal claim, pre-register the gene partition against an external T-helper
regulator database and recompute.

% ===========================================================================
\section{Prioritized roadmap}

Table~\ref{tab:roadmap} orders the reinforcement actions by return on effort. We separate what
is achievable \emph{immediately and in silico} (and would strengthen a hackathon submission this
week) from what requires \emph{new wet-lab data} (and would strengthen a journal submission).

\begin{table}[h]
\centering
\small
\renewcommand{\arraystretch}{1.3}
\begin{tabularx}{\textwidth}{@{}c X l l@{}}
\toprule
\textbf{Priority} & \textbf{Action} & \textbf{Resolves} & \textbf{Status / effort} \\
\midrule
1 & Held-out-modality certificate test (hide activation arm, predict from knockdowns) & L1 & \sev{sev3}{done in silico} (synthetic; awaits dual-arm data) \\
2 & Constraint ablation: NNLS vs unconstrained vs nearest-single & L4 & \sev{sev3}{done in silico} \\
3 & Magnitude-capped recipes with reliability bands & L2 & \sev{sev3}{done in silico} \\
4 & Publish scoring rubric + capability matrix; adversarial axis test & L3 & in silico, days (open) \\
5 & Report per-target reachable-cosine ceiling & L5 & \sev{sev3}{done in silico} \\
6 & Retrospective cross-reference of certificate genes to activation screens/genetics & L1 & \sev{sev3}{done in silico} \\
6b & DEG-magnitude-weighted re-evaluation (signal-dilution stress test) & L5 & \sev{sev3}{done in silico} \\
6c & Cross-cell-type transfer probe (K562/RPE1 essential screens) & L6 (partial) & \sev{sev3}{done in silico} \\
7 & Shared-task comparison vs a forward predictor on a common target & L3 & in silico, weeks \\
8 & In-domain double-knockdown panel in CD4\textsuperscript{+} T cells & L2 & \textbf{wet-lab}, months \\
9 & CRISPRa arm testing the \Th{2}$\rightarrow$\Th{1} certificate & L1 & \textbf{wet-lab}, months \\
10 & Per-donor / multi-system dictionaries for LODO generalization & L6 & \textbf{wet-lab}, months \\
\bottomrule
\end{tabularx}
\caption{Reinforcement actions ordered by return on effort. Actions 2, 3, 5, 6, 6b and 6c have
been \emph{completed} with the data already in hand, and action~1 is verified on synthetic
ground truth and awaits a dual-modality dataset; actions~4 (publish the rubric and run the
adversarial axis test for L3) and~7 (a shared-task comparison) remain open computational work,
and actions 8--10 require new experiments and define the path to a journal-strength result. The
two critical limitations (L1, L2) each retain a definitive wet-lab test (actions~9 and~8) that
in-silico work cannot substitute for, so the critical severities stand despite the completed
mitigations.}
\label{tab:roadmap}
\end{table}

\section{Honest scorecard}

Table~\ref{tab:scorecard} is our own assessment of how well each pillar of the paper is currently
supported. We include it because a hackathon reviewer's time is best spent knowing where we are
confident and where we are not.

\begin{table}[h]
\centering
\small
\renewcommand{\arraystretch}{1.3}
\begin{tabularx}{\textwidth}{@{}l c X@{}}
\toprule
\textbf{Claim} & \textbf{Support} & \textbf{Basis / caveat} \\
\midrule
Reachability is convex-cone membership & \sev{sev3}{strong} & Mathematical; NNLS solves it exactly, KKT residual $\approx\!10^{-11}$. \\
Reachable recipes recover known biology & \sev{sev3}{strong} & GATA3/TBX21 positive controls behave as required; held-out $z$ high. \\
Headline is honest, not overstated & \sev{sev3}{strong} & Held-out (not in-sample) reported; calibration gap disclosed; ceiling argued. \\
Operator transfers across datasets & \sev{sev2}{moderate} & K562 held-out ($z\approx37$) plus a K562/RPE1 probe: \emph{direction} transfers (median $0.35$), the minimal \emph{recipe} does not (Jaccard $0.11$). \\
Additivity supports multi-gene recipes & \sev{sev2}{moderate} & Median cosine $0.71$ (out-of-domain); recipes now magnitude-capped with reliability bands ($0.92$--$0.96$); in-domain doubles still unmeasured. \\
Non-negativity constraint earns its place & \sev{sev3}{strong} & Ablation: constraint costs $0.018$ cosine, saves $\approx3{,}537$ negative weights, beats nearest-single $12/12$; sole source of the certificate. \\
Headline cosine is close to its modality ceiling & \sev{sev3}{strong} & $0.448$ is $71\%$ of the $\sqrt{f_{\mathrm{LOF}}}$ knockdown-only ceiling (atlas mean $61\%$); verdict strengthens under DEG weighting. \\
Disjointness across $92$ methods & \sev{sev2}{moderate} & Structured survey with explicit rubric; not externally re-scored. \\
Certificate ranking is reproducible & \sev{sev3}{strong} & Stable under gene-axis split (cross-half $\rho\approx0.65$, $z\approx35$); top-$30$ retained. \\
Certificate is biologically actionable & \sev{sev1}{weak} & Interpretable, optimal, reproducible, but \emph{never tested at the bench}; top genes are markers/negative regulators, so direction may be backwards. \\
\bottomrule
\end{tabularx}
\caption{Self-assessment of claim support. The framework's mathematical and honesty properties
are strongly supported and have been reinforced by the completed in-silico ablations (constraint,
ceiling, DEG-weighting, split-stability, cross-cell-type); its most novel biological
claim---that infeasibility certificates name genes worth \emph{activating}---remains the least
validated and is the top priority for reinforcement, and the certificate cross-reference shows
its interpretation needs care (the top genes are state-markers and negative regulators, not the
canonical master activators).}
\label{tab:scorecard}
\end{table}

% ===========================================================================
\section{Wet-lab design specifications}
% ===========================================================================
\noindent
The reinforcement items above name three tests that cannot be done in silico on the current
data: a \emph{functional} endpoint (the transcriptome-to-phenotype gap flagged in L1 and by the
pharma review), an \emph{in-domain double-knockdown} panel (L2), and \emph{per-donor} effect
vectors for a leave-one-donor-out test (L6). Each is specified below at the level of detail a
collaborating wet-lab needs to cost and run it, with a pre-registered decision rule so the result
is interpretable however it turns out. Sample sizes are order-of-magnitude estimates for planning,
not a substitute for a statistician's power calculation on the chosen effect size.

\subsection{W1 \quad Functional endpoint: does a reachable recipe move a T\textsubscript{H}1 \emph{phenotype}, not just its transcriptome?}
\sev{sev1}{Resolves the transcriptome$\rightarrow$phenotype gap (L1, pharma item 1). Wet-lab, weeks.}
\begin{itemize}
\item[] \textbf{Objective.} The entire verdict is a transcriptional cosine; a held-out cosine of
$0.45$ says the recipe reproduces a \Th{1} \emph{signature}, not that the cell acquires \Th{1}
\emph{function}. This experiment tests whether the minimal reachable knockdown recipe moves an
orthogonal, functional readout in the predicted direction.
\item[] \textbf{Design.} Primary human CD4\textsuperscript{+} T cells, CRISPRi (dCas9--KRAB), the
same system as the dictionary. Arms: (i)~the oracle's minimal \Th{2}$\rightarrow$\Th{1} recipe
(top reachable knockdowns, $k\!\approx\!5$--$7$); (ii)~a canonical \Th{1} polarizing-cytokine
positive control (IL-12\,+\,anti-IL-4); (iii)~a non-targeting-guide negative control; (iv)~single-gene
knockdown of each recipe member, to show the combination beats its parts. $n\!\ge\!4$ donors,
$\ge\!2$ guides per gene, cells assayed at rest and after restimulation.
\item[] \textbf{Functional readouts (the point of the experiment).} IFN-$\gamma$ secretion
(intracellular stain and ELISA/Luminex on supernatant); T-bet protein (flow); proliferation
(CTV dilution); and, if resourced, an \emph{in-vitro} killing or help assay. Confirm the
transcriptional shift in the same wells (targeted \Th{1}/\Th{2} qPCR panel or mini
RNA-seq) so the transcriptome$\rightarrow$phenotype link is measured, not assumed.
\item[] \textbf{Analysis \& pre-registered decision rule.} Primary endpoint: \%\,IFN-$\gamma^{+}$
in the recipe arm vs.\ non-targeting control, mixed model with donor as random effect.
\emph{Success:} recipe raises the functional endpoint significantly and in the same direction as
the polarizing-cytokine control, at a fraction of the effect that scales with the recipe's
predicted reachable fraction ($\sim\!0.4$). \emph{Informative failure:} transcriptional shift
reproduced but functional endpoint unmoved $\Rightarrow$ the oracle predicts signature but not
function, which is the honest boundary of a signature-matching method and is itself a result.
\item[] \textbf{Scale.} $\sim\!4$ donors $\times\!\sim\!10$ conditions $\times\!2$ guides,
one 96-well CRISPRi campaign with flow + secretion readout; weeks once cells and guides are in hand.
\end{itemize}

\subsection{W2 \quad In-domain double knockdowns: calibrate additivity where the oracle actually runs}
\sev{sev1}{Resolves L2 (additivity borrowed from K562 CRISPRa). Wet-lab, weeks.}
\begin{itemize}
\item[] \textbf{Objective.} Every recipe and certificate assumes co-perturbation effects add. The
only in-hand calibration is $126$ K562 CRISPRa doubles (median cosine $0.71$, magnitude-saturation
ceiling $M^{\star}\!=\!13.9$)---a different cell type \emph{and} a different modality than the
headline CD4\textsuperscript{+} CRISPRi result. This panel measures additivity in-domain.
\item[] \textbf{Design.} $\sim\!60$--$100$ CD4\textsuperscript{+} CRISPRi double knockdowns, plus
both constituent singles and non-targeting controls, read out by pseudobulk or single-cell RNA-seq
(the same effect definition as the dictionary: z-scored log fold-change vs.\ non-targeting pool).
\emph{Pair selection is the design's leverage}: (a)~pairs the oracle actually proposes in
\Th{2}$\rightarrow$\Th{1} and aging recipes; (b)~pairs spanning the combined-magnitude range around
the fitted ceiling $M^{\star}\!=\!13.9$, so the saturation law is tested where it bends, not only
in its flat regime; (c)~a few canonical-regulator pairs (e.g.\ involving \emph{GATA3},
\emph{TBX21}) as interpretable anchors.
\item[] \textbf{Analysis \& pre-registered decision rule.} For each double, cosine and magnitude
ratio of the measured effect vs.\ the sum of singles; regress the additivity deficit on combined
magnitude and test the K562-fit saturation ceiling as an out-of-sample prediction.
\emph{Success:} in-domain median cosine within a pre-set margin (say $\pm0.1$) of the K562 $0.71$
\emph{and} the magnitude-deficit slope reproduced $\Rightarrow$ the borrowed calibration transfers,
and the magnitude-capped recipe reliability bands are validated. \emph{Informative failure:}
in-domain additivity materially worse $\Rightarrow$ recipes above a measured magnitude must carry a
wider band or be capped; this directly sets the recipe-size ceiling.
\item[] \textbf{Scale.} One $\sim\!200$--$300$-guide-arm CRISPRi screen (doubles + singles +
controls, $\times$ replicates) with sequencing readout; the smallest experiment that converts
additivity from a borrowed assumption into an in-domain measurement.
\end{itemize}

\subsection{W3 \quad Per-donor effect vectors and leave-one-donor-out generalization}
\sev{sev2}{Resolves L6 (donor-collapsed dictionary). Wet-lab / re-analysis, weeks--months.}
\begin{itemize}
\item[] \textbf{Objective.} The primary dictionary is collapsed across four donors, so a genuine
leave-one-donor-out (LODO) test---does a recipe built on donors $1$--$3$ still reach the target in
donor $4$?---is impossible on the released data. Pharma needs a nomination to survive patient
genetic background before it is actionable.
\item[] \textbf{Design (two routes).} (i)~\emph{Re-analysis}, if per-donor cells are recoverable
from the original screen: rebuild the effect matrix per donor rather than pooled, giving four
donor-resolved dictionaries at once. (ii)~\emph{New data}, otherwise: repeat the CRISPRi effect
measurement for the recipe-relevant generators (not the full genome) across $\ge\!4$--$6$ donors
spanning HLA and sex, sufficient for a LODO curve on the headline axes.
\item[] \textbf{Analysis \& pre-registered decision rule.} Build the cone on all-but-one donor,
evaluate reachable cosine, held-out cosine, and the recipe membership on the held-out donor; repeat
for each fold. Report the cross-donor mean and spread, and the stability of the top recipe genes
(rank overlap across folds). \emph{Success:} held-out-donor reachable cosine within the
generator-noise band established in silico (Phase-1 dictionary bootstrap) and stable top-recipe
membership $\Rightarrow$ the nomination is genotype-robust. \emph{Informative failure:} recipe
identity swings by donor $\Rightarrow$ the external-validity claim stays ``demonstrated once'' and
donor covariates must enter the model.
\item[] \textbf{Scale.} Route (i) is a re-analysis (no new wet-lab) contingent on data access;
route (ii) is a targeted, generator-restricted CRISPRi campaign across donors. Either yields the
first cross-genetic-background test of the oracle.
\end{itemize}

\noindent\emph{Relation to the in-silico work.} W1--W3 are deliberately the items that
\emph{cannot} be closed at a keyboard. The generator-uncertainty bootstrap, effective-rank/group-%
sparse geometry, held-out-modality certificate cross-reference, directional-genetics check, and
forward-predictor head-to-head (this reinforcement round) tighten everything that can be tightened
in silico; W1 (function), W2 (in-domain additivity), and W3 (cross-donor) are the residual
evidence a drug-development reviewer would still require before acting on a nomination.

\section*{Bottom line}
The manuscript's mathematics and its intellectual honesty are its strongest assets, and its
most novel idea---the constructive infeasibility certificate---is its least validated. The
highest-value work is therefore not more math but a single wet-lab test of a certificate, with
an in-silico held-out-modality test as the immediate proxy. If we could do exactly one more
thing, it would be to hide an activation arm from a dual-modality screen and show that the
certificate, built from knockdowns alone, predicts what activation the screen found to work.
Everything else in this document is a way of making the paper harder to dismiss; that one
experiment is a way of making it right.

\end{document}

